\section{The White Dove}

Says Abdallah al-Bahrani: ``A white dove entered Fatima's garment at the house of Umm Salamah before her meeting with the prophet Muhammad.''\hyperref[_edn1]{ i } As is well known, in traditional Christian (Catholic and Eastern Orthodox) iconography, the white dove is a symbol of the Holy Spirit. In traditional Christian iconography, the Virgin Mary is shown in the presence of a white dove at the time of the Miraculous Conception of Jesus, and at the time of the Annunciation. Because of its connection with the beautiful \emph{Magnificat} prayer, the Annunciation is a common motif of Catholic and Eastern Orthodox iconography. In all icons of the Annunciation, the white dove is shown in the presence of the Virgin Mary. The white dove appeared at the Baptism of Jesus. In all Catholic and Eastern Orthodox icons of the Baptism of Jesus, the white dove is shown hovering over the head of Jesus. In the Gospel According to St. John, it is stated: ``And John (the Baptist) gave testimony saying: `I saw the (Holy) Spirit coming down as a dove from Heaven, and he (the dove) remained upon him (Jesus)'\,'' (John 1:32).

The early Church Fathers wrote a good deal concerning the symbolism of the white dove. White, of course, is symbolic of purity, but there is more; why a white dove and not a white owl or white sea gull? Tertullian wrote: ``The Holy Spirit came in the form of a dove in order that the nature of the Holy Spirit might be made plain by means of a creature of utter simplicity and innocence. For the dove's body has no gall.''\hyperref[_edn2]{ ii }

With his customary eloquence, St. John Chrysostom said:

\begin{quotex}
But why in the form of a dove? The dove is a gentle and pure creature. Since then the (Holy) Spirit, too, is ``a Spirit of gentleness,'' he appears in the form of a dove, reminding us of Noah, to whom, when once a common disaster had overtaken the whole world and humanity was in danger of perishing, the dove appeared as a sign of deliverance from the tempest and, nearing an olive branch, published the good tidings of a serene presence over the whole world.\hyperref[_edn3]{ iii }
\end{quotex}


Here is a selection on Fatima from a discourse attributed to the Prophet Muhammad: ``She (Fatima) shall find herself humiliated after being loved and well treated during the lifetime of her father. Then God will console her with angels who will address her with the words He addressed to Mary, daughter of Imran (Biblical Joachim). They will say to her, `O Fatima, God has chosen thee and purified thee; He has chosen thee above all women.'\,''

Note the close parallel of this with the Catholic Latin prayer ``Ave Maria,'' or ``Hail Mary'':

\begin{quotex}
\emph{Ave Maria, gratia plena, Dominus tecum; benedicta tu in

mulieribus, et benedictus fructus ventri tui Jesus. Sancta Maria, Mater Dei, ora pro nobis pecatoribus nunc et in hora mostis nostrae. Amen.}

Hail Mary, full of grace, the Lord is with thee; blessed art thou among women and blessed is the fruit of thy womb, Jesus. Holy Mary, Mother of God, pray for us sinners now and at the hour of our death. Amen.
\end{quotex}


The above Catholic prayer ``Ave Maria'' (Hail Mary) is based on the Gospel according to St. Luke: ``And the Angel, having come in, said to her: `Hail Mary, full of grace, the Lord is with thee: Blessed art thou among women'\,'' (Luke 1:28). In the Qur'an, it is similarly stated: ``And the angels said: `O Mary! Truly, God has chosen you and purified you and chosen you above the women of the worlds'\,'' (Qur'an 3:42).

The Prophet Muhammad continues:

\begin{quotex}
As for my daughter, Fatima, she is the mistress of the women of the worlds, those that were and those that are to come, and she is part of me. She is the human \emph{houri} who when she enters her prayer chamber before God, exalted be He, her light shines to the angels of Heaven as the stars shine to the inhabitants of the earth. Thus when I saw her, I recalled what will be done to her after me. I could see how humiliation shall enter her home, her sanctity shall be violated, her rights usurped, her inheritance denied and her troubles multiplied. She shall lose her child (through miscarriage), all the while crying out, ``Oh my Muhammad,'' but no one will come to her aid. After me she will remain sorrowful and grieved and weeping, at times recalling the cessation of revelation from her house, at other times my departure from her. When night comes upon her, she shall feel lonely, missing my voice, which she was accustomed to hearing as I recited the Qur'an by night. She shall find herself humiliated after being loved and well treated during the lifetime of her father. Then God will console her with angels, who will address her with the words He addressed to Mary, daughter of Imran. They will say to her, ``O Fatima... God has chosen you, and purified you; He has chosen you above all women. (Fatima)... be obedient to your Lord, prostrating and bowing before Him.''

Then her pains will commence and she will fall, and Mary daughter of Imran will come to her to nurse and console her in her sickness. She shall then say, ``O Lord, I truly despise this life and have become troubled with the people of this world; let me therefore depart to my father.'' Thus she will be the first to come to me from my family. She will come to me sorrowful and heavy with grief, persecuted and martyred. Then will I say, ``O God, curse those who wronged her, punish those who persecuted her, humiliate those who humiliated her, and consign eternally into Your fire him who hit her side so that she lost her child.'' Then the angels will reply: ``Amen.'' ... A Bedouin of one of the tribes in the neighborhood of Medina came to the Prophet (Muhammad) who was sitting with his companions, reviling him and calling him a magician and a liar. He had hidden in his sleeve a small lizard, which he had caught in the desert. He let the animal go and the Prophet called it to him, asking it, ``Do you know who I am?'' The animal answered, ``You are Muhammad, the Apostle of God.'' In astonishment and recognition of the Prophet's claims and forbearance, the Bedouin embraced Islam. But he was poor and hungry and none of the companions had anything to give him to eat. Confident of Fatima's generosity and compassion, the Prophet sent Salman, the Persian, to her seeking food for the hungry man. She had nothing but her own clothes, so she sent her cloak to be pawned with Simon the Jew for a bushel of barley and a tray of dates. She baked the barley, after grinding it with her own hands, and sent the bread and dates to feed the new Muslim. With joy the Prophet came to her, but found her pale with hunger and her two children, Hassan and Hussayn, asleep, trembling like slaughtered birds from hunger as no one in the house of `Ali (Ibn Abi Talib) had tasted anything for three days. The Prophet saw this, and his eyes were filled with tears, and he did not know what to do.

Fatima then entered her chamber and prayed a few \emph{rak'ahs} (daily prayer cycles), after which she invoked God saying, ``O Lord, send to us a banquet \emph{(ma'idah)} from Heaven as You have sent it to the children of Israel. They disbelieved it, yet we will be believers in it. As she finished her prayer, a banquet was sent from Heaven and they all ate. The Prophet, with joy and gratitude, exclaimed, ``Thanks be to God Who has granted me a child like Mary who, ... whenever Zechariah went in to her in the Sanctuary, he found her provisioned. ``Mary,'' he said, ``How does this come to you?'' ``From God'', she said.''\hyperref[_edn4]{ iv }
\end{quotex}


Though generally denounced as idolatry in Sunni Islam as well as Protestantism, intercession (Arabic s\emph{hafa'a}, Persian;
\emph{miyanji gari}) is practiced in Shi`ism as well as Catholicism.

Catholics pray to the Virgin Mary and the saints, and Shi`ahs pray to the Imams and the saints. Ayatollah Murtaza Mutahhari says:

\begin{quotex}
In the ziyarats which we all recite and which we regard as part of Shi`ah doctrine, we say: ``I testify that you see where I stand; you hear what I say and return my salutation.'' It is to be noted that we address that to an Imam who is dead. From our point of view in this respect there is no difference between a dead and a living Imam. It is not that we say so to a dead Imam only. We say: ``Peace on you (Imam) `Ali Ibn Musa al-Riza (8th Imam). I admit and testify that you hear my salutation and return it.\hyperref[_edn5]{ v }
\end{quotex}


Indeed, Ayatullah Murtaza Mutahhari says that denial of intercession and of the immortality of the soul is materialism, and thus very near to pure atheism.\hyperref[_edn6]{ vi } In effect, Ayatullah Mutahhari is saying that any Shi`a who denies intercession is not only a heretic but very near to being an atheist, the very worst sort of unbeliever.

The Sixth Imam, Ja`far al-Sadiq said: ``He (the Imam Hussayn) sees the one who weeps for him, and so he seeks God's forgiveness for him out of compassion, and asks his fathers (`Ali and Muhmmad) to seek pardon for his sins.''\hyperref[_edn7]{ vii } D. K. Crow explains: ``While Hussayn was killed, he nevertheless is alive in Heaven, where alongside the other members of his family, he hears and responds to the prayers of his shi`a, and mediates remission of sins.''\hyperref[_edn8]{ viii }

In \emph{al-Mizan}, his monumental commentary on the Qur'an, Allamah Tabataba'i devotes many pages to a defense of the doctrine of intercession.\hyperref[_edn9]{ ix } As one might expect, the defense of the doctrine of intercession by Ayatollah Mutahhari and Allamah Tabataba'i closely resembles the defense of the same concept by Catholic theologians against attacks against it by Protestants, Modernists and secularists.

The concept of intercession is beautifully expressed by the nineteenth century Persian poet Qa'ani in reference to the Imam Hussayn:

\begin{quotex}
\emph{How was it with him? He fell a martyr! Where? In the Plain of

Karbala!}\\
\emph{When? On the tenth of Muharram! Secretly? No, in

public!... }\\
\emph{Was not the dagger ashamed to cut his throat?}\\

\emph{It was! Why then did it do so? Destiny would not excuse it!}\\

\emph{Wherefore? In order that he (Hussayn) might become an intercessor

for mankind!}\\
\emph{What is the condition of his intercession? Lamentation and

weeping!\hyperref[_edn10]{\textbf{ x }}}
\end{quotex}


Intercession is attributed to the Prophet Muhammad and the Ahl al-Bait by the fourth Shi`a Imam, Imam Zain al-Abidin `Ali Ibn Hussayn, himself son of the Martyr of Karbala and therefore a member of the Ahl al-Bait in his own right.\hyperref[_edn11]{ xi } This demonstrates that the concept of intercession appears very early indeed in the history of Shi`ism.

Intercession is also attributed to Fatima al-Zahra, the \emph{Mater Dolorosa} (Sorrowful Mother) of the martyred Imam Hussayn.\hyperref[_edn12]{ xii } The texts concerning this are far too long and numerous to quote here.

The village of Fatima in Portugal, where the Virgin Mary appeared to shepherd children, is well-known. That the Virgin Mary should appear in a village named for Fatima bint Muhammad Al-Zahra is so astounding a ``coincidence'' that many see in it the hand of Divine Providence, and it is difficult to disagree. Of course, Protestants denounce what they call ``worship of Mary'': as G. K. Chesterton said in his poem ``Don Juan de Austria'': ``And Christian (Protestant) hateth Mary, whom God kissed in Galilee.''

If anyone doubts that Shi`a Islam, and perhaps Sunni Islam, is closer to traditional (Catholic and Eastern Orthodox) Christianity than is Protestantism, compare the attitude of the Shi`as -- and to a lesser extent the Sunnis -- towards Fatima and the Virgin Mary to the Protestant complaint about ``worship of Mary,'' which shows not only the Nestorian basis of Protestantism, but also the Manichaean base.

Moojan Momen states:

\begin{quotex}
Whereas in Sunni Islam there is a direct relationship between the believer and God as revealed in the religion of Islam, in Shi`i Islam there is something of a triangular relationship. While for some things, such as the daily obligatory prayers, the individual is in direct relationship to God, in other matters he looks (usually through the mediation of the local mulla) to the \emph{marja' at-taqlid}, who is regarded as being in a more direct relationship with God. Indeed, in the minds of many of the less educated, the `ulama and the marja' are intermediaries between them and God, and the relationship is not so much triangular as hierarchical.\hyperref[_edn13]{ xiii }
\end{quotex}


The parallel with the relationship of the Catholic or Eastern Orthodox layman to his Church hierarchy and that between the Shi`a layman and the Shi`a hierarchy is obvious. It is sometimes said, erroneously, that Islam has no sacraments. However, it is true that there is nothing in Islam which corresponds exactly with the Catholic and Eastern Orthodox sacraments of Confession and Communion, or Eucharist. However, even this difference is not so great as some imagine, since the concept of doing penance for one's sins is very much a part of Islam. The Protestant principle of \emph{sola fide}, i.e., ``justification by faith alone,'' or ``faith without works,'' is alien and repugnant to Islam. For example, the Qur'an states: ``Truly those who believe and do good works and humble themselves before their Lord, these are the dwellers of Paradise, and they will abide therein forever'' (Qur'an 11:23).

I have heard Protestants speak of the Virgin Mary in terms not far removed from those in the \emph{Talmud, Sanhedrin 51A}, in which the Virgin Mary is referred to as a whore for the Roman Army. Thus the Irish Catholic joke:

\begin{quotex}
A man climbs to the top of a tall building and yells: ``I am going to jump, I am going to commit suicide.'' Someone calls the police, and an Irish policeman is sent to handle the emergency.

``For the love of your father, don't jump'', says the Irish policeman.

``I don't have a father'', says the potential suicide.

``For the love of your mother, don't jump'', says the Irish policeman.

``I don't have a mother either, I am an orphan'', says the potential suicide.

``For the love of the Blessed Virgin don't jump'', says the Irish policeman.

``Who the hell is the Blessed Virgin?'', says the potential suicide.

``Go ahead and jump, you Protestant bastard'', says the Irish policeman.
\end{quotex}


In the Qur'an, the Virgin Mary has a whole Sura named for her. In the Qur'an the Virgin Mary is mentioned 34 times, more than in the gospels. Here are three representative references to the Virgin Mary in the Qur'an: ``And recall, O Our Apostle Muhammad, when the angels said: `O Mary! Truly, God has chosen you and purified you and chosen you above the women of the worlds'\,'' (Qur'an 30:42). ``O Our Apostle Muhammad, remember Mary, who guarded her chastity. We breathed into her Our Spirit, and we made her a sign to all peoples'' (Qur'an 21:91). ``And Mary the daughter of Imran, who guarded her chastity; we breathed into her body Our Spirit, and she testified to the truth of the words of her Lord and His Scriptures, and she was of the obedient ones'' (Qur'an 66:12).

The Shi`a, in particular, have a special reverence for the Virgin Mary, which is linked to their reverence for Fatima, ``the younger Mary.'' If the Virgin Mary is indeed the very touchstone and hallmark of traditional Catholic and Eastern Orthodox truth, there can be no doubt that Shi`a Islam is far closer to traditional Catholic and Eastern Orthodox Christianity than is Protestantism.

In another place we noted that the philosopher Mortimer J. Adler, born into a Jewish family, converted to Christianity because, of the three Abrahamic faiths, i.e., Judaism, Christianity, and Islam, only Christianity affirms God's immanence as well as His transcendence. We also noted that in this Mr. Adler is mistaken, because Islam, save certain aberrant sects, such as Wahhabism and the Taliban, vigorously affirms God's immanence as well as His transcendence.

The Christian monuments of Bethlehem, including the Church of the Nativity, as well as the tomb of St. John the Baptist in Damascus and two places in Jerusalem sacred to the memory of the Virgin Mary, are considered \emph{Ziarats} (places of pilgrimage) by the Shi`as.\hyperref[_edn14]{ xiv } Christian pilgrims to Bethlehem should not be surprised to find themselves joined by Shi`a Muslims.

Nearly everyone has heard of Blaise Pascal (1623 -- 1662), great scientist and Catholic thinker. Particularly famous is ``Pascal's wager'': If one accepts the Catholic faith and is right, then he gains everything, but if he is wrong and the atheists are right, he loses nothing; on the contrary, if one chooses atheism, and is wrong, then he loses everything, but if right he gains nothing.\hyperref[_edn15]{ xv } Pascal had a forerunner who lived nearly one thousand years before him in the figure of Ja`far al-Sadiq, the Sixth Shi`a Imam (702 -- 765). Here is the story: Once during the Hajj in Mecca, the atheist Ibn Abi al-Awja asked Imam al-Sadiq: ``How long will these oxen (Muslims) continue to plow this desert (religion)?'' The Imam responded by citing the Qur'anic verse prohibiting disputatious arguments during the hajj. After the hajj season, the Imam al-Sadiq came to Ibn Awja and said: ``Let me answer with your own style of logic. If there is no God, and everything is absurd, then Muslim worshippers lose nothing by their worship; but if there is a God and Muhammad is His messenger, then woe to you on Judgment Day.''\hyperref[_edn16]{ xvi }

Change the word ``Muslim'' to ``Catholic'', and the Imam al-Sadiq is repeating ``Pascal's wager'' virtually word for word, though nearly one thousand years before the time of Pascal. It is most unlikely, very nearly impossible, that Pascal could have known of the words of the Imam al-Sadiq. Coincidence? Or, perhaps, another example of the affinity between Shi`ism and traditional Catholicism?

All Muslims, whether Sunni or Shi`a, have a great reverence for Jesus, as anyone knows who has read the Qur'an. However, Shi`ites seem to have a particular reverence for Jesus, and, compared to the Sunnis, for St. John the Baptist. Jesus, along with the Twelfth Imam, is called ``the Alive and Awaited.''\hyperref[_edn17]{ xvii }

Among the Shi`as (particularly Iranian Shi`as), Jesus holds a place far superior to that of all the other prophets before Muhammad, as is clearly expressed here:

\begin{quotex}
Seven and a quarter centuries before (Jesus) Christ, Rome was founded, and in the succeeding centuries extended her rule far and wide. Not long after Rome's foundation, Zoroaster arose in Iran and substituted for the magic of Magianism a rational and moral relationship between man and the God of Good in the eternal battle against Evil (i.e., Ahriman or Satan). In almost the same century Confucius and Lao-Tse in China and Gautama the Buddha in India laid the basis of the philosophy which was developed by Socrates, Plato and Aristotle in Greece during the succeeding century. All this found consummation in the birth and life of Jesus Christ, Who proclaimed the call to reform human society, to rescue mankind from the pollutions of Judaistic materialism, to extirpate corruption and internecine combat, and to raise humanity towards ethical and spiritual purification.\hyperref[_edn18]{ xviii }
\end{quotex}


Here Jesus is seen as superior to all prophets before Him, and as heir to Zoroaster as well as Abraham. As we shall see, both the Christian and Zoroastrian traditions affirm this. This is indicated in the accounts in the Gospel According to St. Matthew of the Wise Men, or ``Magi''. The very earliest sources identify the Wise Men as Persians and Zoroastrians. These sources include early apochryphal accounts of the childhood of Jesus, which are considered authoritative and reliable, but are not included in the New Testament Canon because they are incomplete and fragmentary and have little theological importance. These very early accounts leave no doubt as to the Persian and Zoroastrian identity of the Wise Men.

It should be noted that nowhere are the magi given names, nor is it said that they were three in number. The number three is derived from the number of gifts which they brought to the infant Jesus, i.e., gold, frankincense and myrrh. Indeed, the Syrian and Armenian traditions claim that the magi were twelve in number.\hyperref[_edn19]{ xix } The Armenian tradition is particularly significant in this respect, because Armenia was always under strong Persian cultural influence. The Arsacid Dynasty, which long ruled Armenia, was of Parthian origin, and for a long time there were many Zoroastrians in Armenia.

It is from these extra-canonical accounts that it is known that the parents of the Virgin Mary were St. Joachim (the Qur'anic \emph{Imran}) and Ste. Anne. The First Gospel on the Infancy of Jesus Christ is certainly of great antiquity, as the earliest references to it so far discovered are from the early Second Century. Various Church Fathers accepted the reliability of this Gospel, including St. Athanasius and St. John Chrysostom.\hyperref[_edn20]{ xx } It was not included among the canonical gospels because, dealing only with the infancy of Jesus, it is fragmentary and has little of theological importance. Here is what it has to say concerning the Wise Men:

\begin{quotex}
And it came to pass, when Lord Jesus was born at Bethlehem, a city of Judaea, in the time of Herod the King; the wise men came from the East to Jerusalem, according to the prophecy of Zoroaster, and brought with them offerings: namely, gold, frankincense and myrrh, and worshipped him, and offered him their gifts... .And having, according to the custom of their country, made a fire, they worshipped it.\hyperref[_edn21]{ xxi }
\end{quotex}


According to legend, no doubt of Zoroastrian origin, the Island of Kuh-i-Khwaja in Lake Helmand in Eastern Iran, called
\emph{Daryacheh-ye-Sistan}, is the site of the castle which was the home

of the Wise Men.\hyperref[_edn22]{ xxii }

Christian tradition from earliest times has always portrayed the Wise Men in Persian garb. So the Christian Tradition has always affirmed that Jesus was heir to Zoroaster as well as to the Old Testament prophets. In Iran much is made of this by Christians, Zoroastrians and Muslims. Remember, according to Islam, Muhammad was the last in a line of prophets which includes Jesus. So, if Jesus was heir of Zoroaster as well as the Old Testament prophets, so, by extension, was Muhammad. Among Muslims, Shi`ahs make much more of this than do Sunnis, in part because Shi`ahs revere Jesus and the Virgin Mary even more than do Sunnis, and because the Imam Hussayn was married to Shahrbanu, a Persian princess born a Zoroastrian. Shahrbanu was therefore the female ancestor of nine of the twelve Imams. In Iran, Zoroastrians call the Imam Hussayn ``son-in-law,'' which is \emph{damad} in Persian, or
\emph{damad-emahbub}, ``beloved son-in-law''. Since \emph{damad} simply

means ``in-law'' or ``relative by marriage,'' a more formal and precise way of expressing it would be \emph{mard ke shauhar-e-dok-}
\emph{htar-e-man ast,} i.e., ``the man who is the husband of our

daughter'' or \emph{mard ke shauhar-e-dokhtar-e-man mahbub ast}, i.e., ``the man who is the beloved husband of our daughter.''

The implications of Jesus being the heir of Zoroaster as well as the Old Testament prophets are enormous, and it is strange that, outside of Iran, not much has been made of it. The above certainly affects the whole concept of the ``Chosen People'' and the idea that ``salvation is (exclusively) from the Jews.'' Iranians are not only not Jewish, they are not even Semites, being Indo-Europeans or Aryans. Note that the name ``Iran'' comes from the same Indo-European or Sanskrit root as ``Arya'' or ``Aryan,'' and the Old Gaelic ``Erinn,'' Modern Gaelic ``Erin,'' the native Celtic name of Ireland.

Shi`as also have a particular reverence for St. John the Baptist. Many Shi`a traditions speak reverently of St. John the Baptist, such as this account:

\begin{quotex}
When asked about the inner meaning of the letters used, it was said that God willed Gabriel to make Zachariah (father of St. John the Baptist) know the sacred names of the Holy Prophet (Muhammad) along with those of the Ahl al-Bait (the House of Muhammad, in other words, Fatima and the Twelve Imams). When Gabriel mentioned the names, Muhammad, `Ali, Fatima and Hussayn, Zachariah felt a great joy and consolation, but with the mention of the name Hussayn, he was filled with grief and sorrow. When he asked the angels about the uncontrollable sorrow he felt, Gabriel informed him of the heart-rending tragedy of Karbala, for the experiences of Hussayn are quite similar to what Yahya (St. John the Baptist) faced. Zachariah was informed of the story of the untold miseries and tortures which Hussayn would suffer and the brutal massacre he would meet. It is said that in conveying the tidings about the tragedy of Karbala, the letter symbols used in the start of this chapter were: (1.) ``Kaf'' for Karbala, where Hussayn was martyred along with the band of his faithful devotees. (2.) ``Ha'', for halakat, the annihilation of the Holy Family. (3.) ``Ya'' for ``Yazid,'' the son of Mu'awiya, who caused the heartless massacre. (4.) ``\,`Ain'' for `\emph{atash}, the killing thirst (along with hunger) which Hussayn and his followers suffered before they were butchered. (5.) ``Saad'' for
\emph{sabr}, the marvelous patience with which Hussayn and his comrades

suffered torture before they drank the cup of martyrdom.

The events connected with the life of the Third Imam Hussayn, the second son of `Ali and Fatima, the second grandson of the Holy prophet, are identical with those of the apostle Yahya (St. John the Baptist). Hussayn, during his journey to Karbala, frequently remembered the apostle Yahya. Like Yahya, Hussayn was also born in six months. Hussayn's name was also peculiar to him as was the name of Yahya (i.e., refers to first use of said names). Hussayn was martyred for opposing the brute Yazid's devilish life as was Yahya for declaring that what the king did was wrong. With the glad tidings of getting a son, Zachariah was informed of what the name of the son was to be. No one before him had been given the name Yahya (John), derived from \emph{hayat}, i.e., life. It is said that the name Yahya was given to him because he was born of a barren woman who, for never being productive and also for having crossed the age of fertility, was virtually dead to womanhood. The issue or the son promised was Yahya (St. John the Baptist), the forerunner of Jesus... .

The statement ``We did not make anyone like him'' means that Yahya was made such that he neither sinned nor had he any inclination towards any sin. He was not even inclined to marry any woman, and he also did not marry. Secondly, no one else was born of a barren woman who was at the advanced age of about eighty years. God's Almighty will manifested once in bringing forth a child (Jesus) from a virgin (Mary) and once in bringing forth a child (St. John the Baptist) from a barren woman at the advanced age of eighty. And never did the heavens mourn for anyone else as much as for Yahya and the Third Imam Hussayn Ibn `Ali, the King of Martyrs.

The Sixth Imam, Ja`far Ibn Muhammad Al-Sadiq, says that the cases of Yahya and Hussayn are similar because: (1) No one had his name before, save Yahya and Hussayn. (2) For no one else did the heavens weep for forty days save Yahya and Hussayn. (3) The murderers of both Yahya and Hussayn were of illegitimate birth. (4) Sufyan Ibn A`seneh narrates from `Ali Ibn Zaid, who got it from his father `Ali Ibn al-Hussayn, that he said, ``When we set out with Hussayn for Kufa, we did not halt at any place but Hussayn mentioned Yahya Ibn Zachariah. One day he said: `The proof of the worthlessness of this world in the view of God is that the head of Yahya Ibn Zachariah was presented to one of the prostitutes of the Israelites.'\,''
\end{quotex}


Similarly, Hussayn's head was presented to the sons of the prostitutes. He set out from Medina knowing well the fate that awaited him and his family. His aim was that his sacrifice should not take place unnoticed and go to waste ineffective, and he succeeded in his divinely planned mission in laying down his all for the sake of truth in such a way that it shook the very throne of the Tyrant and the heart of every Muslim in particular and humanity in general.\hyperref[_edn23]{ xxiii }

In \emph{Al-Mizan}, his monumental commentary on the Qur'an, Allamah Tabataba'i says:

\begin{quotex}
Allah gave him (Zechariah) a son, Yahya, the prophet most similar to Isa (Jesus) (peace be on both). He was given all the qualities of perfection and excellence which Isa and his truthful mother, Maryam (the Virgin Mary) were granted. It was for this reason that Allah named him Yahya and sent him to verify a word from Allah, and made him honorable and chaste as well as a prophet, from among the good ones. It was the nearest that any man could resemble Maryam and her son Isa, peace be on them all... .Yahya was made to resemble Isa as completely and perfectly as was possible. Yet Isa had precedence over Yahya because his creation and birth were firmly decreed long before the prayer of Zakariyya for Yahya was accepted. That is why Isa was given superiority over Yahya, and made an \emph{ulu l-azam} apostle, bringing a new shari`ah (law) and a new book (the Gospel). Apart from some necessary dissimilarities, Yahya and Isa resembled each other to the maximum extent possible.\hyperref[_edn24]{ xxiv }
\end{quotex}

\hyperref[_ednref1]{ i } Ayoub,

op. cit., p. 72, citing \emph{Maqtal al-Awalim}, volume XVI of
\emph{Awalim al-`Ulum} (Tabriz, Iran, undated), p. 4.


\hyperref[_ednref2]{ ii }

\emph{Ancient Christian Commentary on Scripture: New Testament II;

Mark}, edited by Thomas C. Oden \& Christopher A. Hall, (Downers Grove, Illinois, 1998), p. 13.

\hyperref[_ednref3]{ iii } Idem.


\hyperref[_ednref4]{ iv } Abu

Ja`far Muhammad ibn babaway al-Qummi, Al-Amali wa al-Majlis (Qum, Iran, 1951), p. 70. Cited in Ayoub, op. cit., pp. 238-240. 31

\hyperref[_ednref5]{ v }

Ayatollah Murtaza Mutahhari, \emph{Man and Universe} (Karachi, Pakistan, 1991), p. 463.

\hyperref[_ednref6]{ vi }

Ayatullah Murtaza Mutahhari, Fundamentals of Islamic Thought (Berkeley, California, 1985), pp. 108-111.

\hyperref[_ednref7]{ vii } Ibn

Qawlawayhi, \emph{Kamil al-Zitadi}, cited in Douglas Karim Crow ``The Death of al-Husayn and the Imamate,'' \emph{Al-Serat} (London, July 1984), p. 93.

\hyperref[_ednref8]{ viii }

Ibid., p. 109.

\hyperref[_ednref9]{ ix }

Allamah Muhammad Husayn Tabataba'i, \emph{Al-Mizan, An Exegesis of the Qur`an}, Volume I (Tehran, 1983), pp. 221-262.

\hyperref[_ednref10]{ x }

Edward G. Browne, \emph{Literary History of Persia}, Volume IV (Cambridge, England, 1959), pp. 178-181.

\hyperref[_ednref11]{ xi } Zayn

al-Abidin `Ali ibn Husayn, \emph{The Psalms of Islam: Al-Sahifat al-Kamilat al-Sajjadiyya}, edited and translated by William C. Chittick (London, 1988), pp. 22, 89. 189, 228-229, 250.

\hyperref[_ednref12]{ xii }

Ordoni, op. cit., pp. 268-288; Tabataba'i, Al-Mizan, \emph{An Exegesis of the Qur'an}, volume I, op. cit., pp. 256-257.

\hyperref[_ednref13]{ xiii }

Moojan Momen, \emph{An Introduction to Shi'i Islam} (New Haven, Connecticut, 1985), p. 235.

\hyperref[_ednref14]{ xiv }

Yousuf N. Lalljee, \emph{Know Your Islam} (Qum, Iran, undated), pp. 214-216.

\hyperref[_ednref15]{ xv }

Blaise Pascal, \emph{Pensees}, Introduction and notes by Charles-Marc des Granges (Paris, 1961), Pensee no. 233, pp. 134-138.

\hyperref[_ednref16]{ xvi }

Michael M. J. Fischer and Mehdi Abedi, \emph{Debating Muslims} (Madison, Wisconsin, 1990), p. 500, note 7. Fischer and Abedi do not give a source.

\hyperref[_ednref17]{ xvii }

\emph{The Holy Qur'an}, Text, Translation and Commentary by S. V. Mir

Ahmed `Ali, Introduction by Agha Haji Mirza Mahdi Pooya Yazdi, op. cit., p. 153a.

\hyperref[_ednref18]{ xviii }

Sayid Mujtaba Rukni Musawi Lari, \emph{Western Civilization Through Muslim Eyes} (Qum, Iran, undated), p. 3.

\hyperref[_ednref19]{ xix }

\emph{Oxford Dictionary of Byzantium}, Alexander P. Kashdan,

editor-in-chief, essay ``Adoration of the Magi'' (Oxford, England, 1991), vol. 1, p. 22.

\hyperref[_ednref20]{ xx }

\emph{The Lost Books of the Bible} (New York, 1979), p. 38. ``Lost'' is

an unfortunate adjective in this case, because the books in this particular anthology, though not included in the New Testament Canon, were never ``lost''. There are several collections of New Testament Apocrypha and/or non-canonical gospels available.

\hyperref[_ednref21]{ xxi }

Ibid., p. 40.

\hyperref[_ednref22]{ xxii }

Richard N. Frye, \emph{The Heritage of Persia} (Cleveland, Ohio, 1963), p. 178.

\hyperref[_ednref23]{ xxiii }

Commentary at bottom of pages of \emph{The Holy Qur'an}: Text, Translation and Commentary by S. V. Mir Ahmed Ali, op. cit., pp. 942 -- 954.

\hyperref[_ednref24]{ xxiv }

Allamah Sayyid Muhammad Husayn Tabataba'i, \emph{Al-Mizan} (Tehran, 1973), volume 5, pp. 261-262.


\flrightit{Posted on 2016-09-13 by Michael McClain}

\begin{center}* * *\end{center}

\begin{footnotesize}
\begin{sffamily}
\texttt{Ulysses on 2016-09-14 at 14:22 said:}

Is this opinion about the infancy gospel common among catholics?

\hfill

\texttt{Cologero on 2016-09-14 at 16:38 said:}

Ulysses, this is what Pope Benedict XVI wrote about it:
\textit{Revelation to the Wise}\footnote{\url{http://www.gornahoor.net/?p=7776}}.


\hfill

\texttt{Mark Citadel on 2016-09-15 at 07:48 said:}

Fascinating read, in all three parts. I always love stuff involving symbology in particular. My article on Shia geopolitics may provide further interest for those wanting to learn how they apply their religious history to contemporary foreign affairs.

\url{http://www.socialmatter.net/2016/09/06/shia-geopolitics-fortress-state/}



\hfill

\end{sffamily}
\end{footnotesize}

